% Generated draft: requires mathematical and source review.
\documentclass[11pt]{article}
\usepackage[margin=25mm]{geometry}
\usepackage{fontspec}
\setmainfont{Latin Modern Roman}
\usepackage{amsmath,amssymb,mathtools,unicode-math}
\setmathfont{Latin Modern Math}
\usepackage{enumitem,needspace,xurl,hyperref}
\setlength{\parindent}{0pt}
\setlength{\parskip}{5pt}
\setlength{\emergencystretch}{3em}
\newcommand{\N}{\mathbb{N}}
\newcommand{\Z}{\mathbb{Z}}
\newcommand{\Q}{\mathbb{Q}}
\newcommand{\R}{\mathbb{R}}
\newcommand{\C}{\mathbb{C}}
\newcommand{\F}{\mathbb{F}}
\newcommand{\A}{\mathbb{A}}
\newcommand{\PP}{\mathbb{P}}
\newcommand{\E}{\mathbb{E}}
\newcommand{\eps}{\varepsilon}
\newcommand{\dd}{\,\mathrm{d}}
\begin{document}
\section*{UnsolvedMath ID: 10000060}
\textbf{Source problem code:} catalogue entry
\subsection{Definitions and mathematical statement}
Let \(G\) range over finite connected vertex-transitive graphs with at least two vertices, and let \(d_G\) be unweighted shortest-path distance on \(V(G)\). Vertex transitivity means that \(\operatorname{Aut}(G)\) acts transitively on vertices; no Cayley, planar, degree, or size restriction is added. Normalize
\[
\widehat d_G=d_G/\operatorname{diam}(G).
\]
For compact metric spaces \(A,B\), their Gromov--Hausdorff distance is the infimum of Hausdorff distances between isometric copies in a common metric space.

For the intended fixed round-sphere metric \(d_S\), set
\[
\delta_* = \inf_G d_{\mathrm{GH}}((V(G),\widehat d_G),(S^2,d_S)).
\]
The question asks which graph attains this infimum, and in particular whether the truncated-icosahedron graph \(H\) does:
\[
d_{\mathrm{GH}}((V(H),\widehat d_H),(S^2,d_S))=\delta_*.
\]
Under the vertex-set convention \(H\) is the usual 60-vertex one-skeleton. Existence or uniqueness of a minimizer is not assumed by the infimum definition.

The sphere's scale and metric need specification. The unit-radius intrinsic round metric has diameter \(\pi\); a diameter-one intrinsic sphere uses \(d_{\mathrm{geo}}/\pi\). The chord metric, even rescaled to diameter one, is another metric. The source does not say which normalization is paired with the diameter-one graphs. It also does not explicitly distinguish the vertex-set graph metric from the continuous length metric on all edge interiors. These affect the objective and cannot be silently interchanged.

The supplied research summary reports no resolution and mentions vertex-transitive approximation work, without an extremal GH-distance theorem. Approximation results alone do not identify this global minimizer.
\subsection{Short English statement}
Which diameter-normalized finite vertex-transitive graph is closest to the intended round sphere in Gromov--Hausdorff distance? Is it the truncated icosahedron? Fix the sphere/graph metric conventions and do not assume the infimum is attained.
\subsection{Sources}
\begin{itemize}
\item Local UnsolvedMath registry, permanent ID \(10000060\), exact source code \texttt{catalogue entry}, ``Closest finite vertex-transitive graph to the round sphere.'' Proposer Itai Benjamini, 2012.
\item Benjamini, \emph{Euclidean vs. Graph Metric}, Question 5.3, recorded PDF page 6; \url{https://www.wisdom.weizmann.ac.il/~itai/erd100.pdf#page=6}. Recorded access 2026-07-29; extracted text checked against rendered source.
\item Recorded preservation copy: \url{https://arquivo.pt/noFrame/replay/20201231041538id_/http://www.wisdom.weizmann.ac.il/~itai/erd100.pdf}.
\item Supplied OPEN-TRIAGE summary mentions Benjamini--Finucane--Tessera approximation theory without a complete citation or theorem resolving this minimization problem.
\end{itemize}
\end{document}
